\begin{tikzpicture}[
    node distance=0.5cm,
    % Styl dla głównych warstw (duże kontenery)
    layer/.style={
        rectangle,
        draw=gray!40,
        thick,
        fill=white,
        rounded corners,
        inner sep=10pt,
        align=center
    },
    % Styl dla etykiet warstw (po lewej stronie)
    layer_label/.style={
        rotate=90,
        anchor=center,
        font=\sffamily\bfseries\small,
        color=gray!70
    },
    % Styl dla komponentów wewnątrz (aplikacje, vm)
    component/.style={
        rectangle,
        draw=black!60,
        thick,
        fill=white,
        minimum height=1.2cm,
        minimum width=2.5cm,
        text centered,
        rounded corners=2pt,
        font=\footnotesize\sffamily,
        drop shadow
    },
    % Strzałki relacji "Runs on"
    runs_on/.style={
        draw,
        -latex,
        thick,
        gray!50,
        dashed
    }
]

    % --- 1. WARSTWA APLIKACJI (GÓRA) ---
    
    % Grupa: System Sterowania (Management)
    \node[component, fill=red!10, text width=2.2cm] (tuner) {
        \textbf{Smart Tuner}\\
        (MAPE-K)
    };
    
    \node[component, right=0.5cm of tuner, fill=red!5, text width=2.2cm] (scalers) {
        \textbf{Scalers}\\
        (HPA, VPA)
    };

    % Grupa: Monitoring
    \node[component, right=0.5cm of scalers, fill=blue!10, text width=2.2cm] (monitor) {
        \textbf{Monitoring}\\
        (Prometheus)
    };

    % Grupa: Workload
    \node[component, right=0.5cm of monitor, fill=green!10, text width=2.2cm] (app) {
        \textbf{Workload}\\
        (FastAPI)
    };

    % Kontener warstwy aplikacji
    \node[layer, fit=(tuner)(scalers)(monitor)(app), label={[anchor=north west, inner sep=5pt]north west:\textit{Warstwa Aplikacji}}] (app_layer) {};
    \node[layer_label, left=0.2cm of app_layer] {Workloads};


    % --- 2. WARSTWA ORKIESTRACJI (K3S) ---
    
    \node[layer, below=0.8cm of app_layer, minimum width=11.5cm, fill=yellow!10, label={[anchor=north west]north west:\textit{Warstwa Orkiestracji}}] (k3s_layer) {
        \vspace{0.3cm}
        \textbf{K3s Kubernetes Cluster} (v1.33)\\
        \footnotesize{CNI: Flannel \quad | \quad CSI: Local Path}
    };
    \node[layer_label, left=0.2cm of k3s_layer] {K3s};


    % --- 3. WARSTWA MASZYN WIRTUALNYCH ---
    
    % Definiujemy 3 VM-ki pod spodem
    
    % POPRAWKA: Zmieniono 'width' na 'text width'
    \node[component, below=1.0cm of k3s_layer, xshift=-3.5cm, fill=blue!5, text width=3cm] (vm1) {
        \textbf{VM-01}\\
        (Control Plane)
    };
    
    % POPRAWKA: Zmieniono 'width' na 'text width'
    \node[component, right=0.8cm of vm1, fill=blue!5, text width=3cm] (vm2) {
        \textbf{VM-02}\\
        (Worker)
    };
    
    % POPRAWKA: Zmieniono 'width' na 'text width'
    \node[component, right=0.8cm of vm2, fill=blue!5, text width=3cm] (vm3) {
        \textbf{VM-03}\\
        (Worker)
    };

    % Kontener warstwy VM
    \node[layer, fit=(vm1)(vm2)(vm3), label={[anchor=north west]north west:\textit{Warstwa Wirtualizacji}}] (vm_layer) {};
    \node[layer_label, left=0.2cm of vm_layer] {VMs};


    % --- 4. WARSTWA SPRZĘTOWA (PROXMOX) ---
    
    \node[layer, below=0.6cm of vm_layer, minimum width=11.5cm, fill=gray!10, label={[anchor=north west]north west:\textit{Warstwa Infrastruktury}}] (proxmox_layer) {
        \vspace{0.2cm}
        \textbf{Proxmox VE Hypervisor}\\
        \footnotesize{Zasoby Fizyczne: CPU, RAM, SSD}
    };
    \node[layer_label, left=0.2cm of proxmox_layer] {Hardware};


    % --- STRZAŁKI RELACJI ---
    
    \draw[runs_on] (app_layer.south) -- node[right, font=\tiny] {Hostowane w} (k3s_layer.north);
    
    \draw[runs_on] (k3s_layer.south -| vm1.north) -- (vm1.north);
    \draw[runs_on] (k3s_layer.south -| vm2.north) -- (vm2.north);
    \draw[runs_on] (k3s_layer.south -| vm3.north) -- (vm3.north);

    \draw[runs_on] (vm_layer.south) -- node[right, font=\tiny] {Wirtualizacja} (proxmox_layer.north);

\end{tikzpicture}
